%% kapa: exam=1 round=22 year=2011 period=1 form=A subject=civil_law subject_ko=민법 seq=32 number=32
%% source: Q-Net 감정평가사 기출문제 (cst00309) artlSeq=5036790 file=extracted/1008/5036790/2011년 제22회 감정평가사 1차시험 기출문제/1교시.hwp
%% conversion: kordoc 4.24.51 -> tools/gen1.py (auto); review=proofread-sonnet; flags=none
\begin{question}{32}
乙이 甲소유의 부동산을 매수한 후 다시 丙에게 매도하고 인도하였는데, 등기는 아직 甲명의로 남아 있다. 다음 설명으로 옳지 않은 것은? (다툼이 있으면 판례에 의함)
\begin{choices}
  \item 甲ㆍ乙ㆍ丙 전원이 중간생략등기의 합의를 한 경우, 丙은 특별한 사정이 없는 한 직접 甲에게 소유권이전등기를 청구할 수 있다.
  \item 甲ㆍ乙ㆍ丙 사이의 중간생략등기의 합의 없이 丙명의로 등기가 된 경우, 중간생략등기의 합의가 없었다는 이유만으로 그 등기를 무효라고 할 수 없다.
  \item 甲ㆍ乙ㆍ丙 사이의 중간생략등기의 합의가 있더라도 甲은 乙이 매매대금을 지급하지 않았음을 이유로 丙의 소유권이전등기청구를 거절할 수 있다.
  \item 만약 乙이 甲에 대한 소유권이전등기청구권을 丙에게 양도하고 이를 甲에게 통지하였다면, 丙은 특별한 사정이 없는 한 甲에게 직접 소유권이전등기를 청구할 수 있다.
  \item 만약 甲소유의 부동산이 토지거래허가구역 내의 토지라면, 甲ㆍ乙ㆍ丙 전원이 중간생략등기의 합의를 하고 甲ㆍ丙을 매매당사자로 하는 토지거래허가를 받아 丙명의의 소유권이전등기를 해 주었더라도 그 등기는 무효이다.
\end{choices}
\end{question}
\answer{4}
